\par\Needspace{4\baselineskip}
\begin{center}
\hypertarget{native-047-article032}{}\pdfbookmark[2]{32.}{article-032}32.
\end{center}
Wir wenden diese Grundsätze nun auf die nützlichsten der oben erläuterten Rechenverfahren an.

I. Verwendet man zur Berechnung der wahren Anomalie aus der exzentrischen Anomalie bei elliptischer Bewegung die Formel VII aus Artikel 8 und setzt man $\varphi$ und $E$ als genau bekannt voraus, so kann in \SourceLeaf{048}$\log\operatorname{tang}(45^0-\frac{1}{2}\varphi)$ und $\log\operatorname{tang}\frac{1}{2}E$ jeweils ein Fehler von $\omega$ auftreten, in ihrer Differenz $=\log\operatorname{tang}\frac{1}{2}v$ folglich ein Fehler von $2\omega$. Der größte Fehler bei der Bestimmung des Winkels $\frac{1}{2}v$ ist daher $\frac{3\omega\,\mathrm{d}\frac{1}{2}v}{\mathrm{d}\log\operatorname{tang}\frac{1}{2}v}=\frac{3\omega\sin v}{2\lambda}$, wobei $\lambda$ den Modul der für diese Rechnung verwendeten Logarithmen bezeichnet. Für den möglichen Fehler der wahren Anomalie $v$ ergibt sich also, in Bogensekunden ausgedrückt, $=\frac{3\omega\sin v}{\lambda}206265^{\prime\prime}=0^{\prime\prime}{,}0712\sin v$, wenn man \textsc{Briggs}sche Logarithmen mit sieben Dezimalstellen verwendet, sodass wir stets bis auf $0^{\prime\prime}{,}07$ Gewissheit über den Wert von $v$ haben können. Bei Verwendung kleinerer Tafeln mit nur fünf Stellen könnte der Fehler bis auf $7^{\prime\prime}{,}12$ anwachsen.

II. Berechnet man $e\cos E$ mithilfe von Logarithmen, so kann der Fehler bis zu $\frac{3\omega e\cos E}{\lambda}$ betragen. Demselben Fehler ist daher die Größe $1-e\cos E$ oder $\frac{r}{a}$ unterworfen. Bei der Berechnung des Logarithmus dieser Größe kann der Fehler somit bis auf $(1+\delta)\omega$ anwachsen, wobei $\delta$ den Betrag von $\frac{3e\cos E}{1-e\cos E}$ bezeichnet. Bis zu derselben Grenze $(1+\delta)\omega$ kann der Fehler in $\log r$ anwachsen, sofern $\log a$ als genau gegeben vorausgesetzt wird. Bei kleiner Exzentrizität bleibt die Größe $\delta$ stets in engen Grenzen. Unterscheidet sich dagegen $e$ nur wenig von 1, so bleibt $1-e\cos E$ sehr klein, solange $E$ klein ist. Dann kann $\delta$ also einen beträchtlichen Wert erreichen, weshalb in diesem Fall die Formel III aus Artikel 8 weniger geeignet wäre. Die Größe $\delta$ lässt sich auch als $\frac{3(a-r)}{r}=\frac{3e(\cos v+e)}{1-ee}$ ausdrücken. Diese Formel zeigt noch deutlicher, wann man den Fehler $(1+\delta)\omega$ vernachlässigen darf.

III. Verwendet man die Formel X aus Artikel 8 zur Berechnung der wahren aus der exzentrischen Anomalie, so ist $\log\BareRoot{\frac{a}{r}}$ einem Fehler von $(\frac{1}{2}+\frac{1}{2}\delta)\omega$ unterworfen und folglich $\log\sin\frac{1}{2}\varphi\sin E\BareRoot{\frac{a}{r}}$ einem Fehler von $(\frac{5}{2}+\frac{1}{2}\delta)\omega$. Daraus ergibt sich als größter möglicher Fehler bei der Bestimmung des Winkels $v-E$ oder $v$ der Wert $=\frac{\omega}{\lambda}(7+\delta)\operatorname{tang}\frac{1}{2}(v-E)$ oder, in Bogensekunden ausgedrückt, wenn man sieben Dezimalstellen verwendet, $=(0^{\prime\prime}{,}166+0^{\prime\prime}{,}024\,\delta)\operatorname{tang}\frac{1}{2}(v-E)$. Bei mäßiger Exzentrizität sind $\delta$ und $\operatorname{tang}\frac{1}{2}(v-E)$ kleine Größen, weshalb dieses Verfahren eine höhere Genauigkeit erlaubt als das in I betrachtete. Letzteres ist dagegen vorzuziehen, wenn die Exzentrizität sehr groß ist und sich der Eins nähert, da $\delta$ und $\operatorname{tang}\frac{1}{2}(v-E)$ dann sehr beträchtliche Werte annehmen können. Mit unseren Formeln lässt sich stets leicht entscheiden, welches Verfahren vorzuziehen ist.

IV. Bestimmt man die mittlere Anomalie aus der exzentrischen Anomalie nach Formel XII aus Artikel 8, so kann der Fehler der mithilfe von Logarithmen berechneten Größe $e\sin E$ und damit auch der Anomalie $M$ selbst bis auf $\frac{3\omega e\sin E}{\lambda}$ anwachsen. Diese Fehlergrenze ist,
\SourceLeaf{049}wenn man sie in Bogensekunden angeben will, mit $206265^{\prime\prime}$ zu multiplizieren. Daraus folgt leicht, dass bei der umgekehrten Aufgabe, bei der man $E$ aus $M$ durch wiederholtes Probieren bestimmt, $E$ um $\frac{3\omega e\sin E}{\lambda}\cdot\frac{\mathrm{d}E}{\mathrm{d}M}\cdot206265^{\prime\prime}=\frac{3\omega ea\sin E}{\lambda r}\cdot206265^{\prime\prime}$ vom richtigen Wert abweichen kann, selbst wenn die Gleichung $E-e\sin E=M$ mit der vollen Genauigkeit erfüllt ist, die die Tafeln zulassen.

Die aus der mittleren Anomalie berechnete wahre Anomalie kann also aus zwei Gründen fehlerhaft sein, wenn wir die mittlere als genau gegeben ansehen. Zum einen kann bei der Berechnung von $v$ selbst aus $E$ ein Fehler auftreten, der, wie wir gesehen haben, stets von geringer Bedeutung ist. Zum anderen kann der Wert der exzentrischen Anomalie selbst bereits fehlerhaft sein. Die Wirkung der zweiten Ursache ergibt sich aus dem Produkt des Fehlers in $E$ mit $\frac{\mathrm{d}v}{\mathrm{d}E}$. Dieses Produkt ist $=\frac{3\omega e\sin E}{\lambda}\cdot\frac{\mathrm{d}v}{\mathrm{d}M}\cdot206265^{\prime\prime}=\frac{3\omega ea\sin v}{\lambda r}\cdot206265^{\prime\prime}=\left(\frac{e\sin v+\frac{1}{2}ee\sin 2v}{1-ee}\right)0^{\prime\prime}{,}0712$, wenn man sieben Stellen verwendet. Dieser Fehler ist bei kleinen Werten von $e$ stets mäßig, kann aber sehr groß werden, wenn $e$ sich nur wenig von der Eins unterscheidet. Die folgende Tabelle zeigt dies, indem sie für einige Werte von $e$ den größten Wert dieses Ausdrucks angibt.

\begin{center}
\begin{tabular}{c|c||c|c||c|c}
$e$ & größter Fehler & $e$ & größter Fehler & $e$ & größter Fehler \\
0,90 & $0^{\prime\prime}{,}42$ & 0,94 & $0^{\prime\prime}{,}73$ & 0,98 & $2^{\prime\prime}{,}28$ \\
0,91 & $0{,}48$ & 0,95 & $0{,}89$ & 0,99 & $4{,}59$ \\
0,92 & $0{,}54$ & 0,96 & $1{,}12$ & 0,999 & $46{,}23$ \\
0,93 & $0{,}62$ & 0,97 & $1{,}50$ & & \\
\end{tabular}
\end{center}

V. Bestimmt man bei hyperbolischer Bewegung $v$ nach Formel III aus Artikel 21 aus den genau bekannten Größen $F$ und $\psi$, so kann der Fehler bis auf $\frac{3\omega\sin v}{\lambda}\cdot206265^{\prime\prime}$ anwachsen. Verwendet man zur Berechnung dagegen die Formel $\operatorname{tang}\frac{1}{2}v=\frac{(u-1)\operatorname{cotang}\frac{1}{2}\psi}{u+1}$ mit den genau bekannten Größen $u$ und $\psi$, so ist die Fehlergrenze um ein Drittel größer, nämlich $=\frac{4\omega\sin v}{\lambda}\cdot206265^{\prime\prime}=0^{\prime\prime}{,}09\sin v$ bei sieben Stellen.

VI. Berechnet man die Größe $\frac{\lambda kt}{b^{\frac{3}{2}}}=N$ nach Formel XI aus Artikel 22 mithilfe von \textsc{Briggs}schen Logarithmen und setzt man dabei $e$ und $u$ oder $e$ und $F$ als genau bekannt voraus, so ist der erste Teil einem Fehler von $\frac{5(uu-1)e\omega}{2u}$ unterworfen, wenn man ihn in der Form $\frac{\lambda e(u-1)(u+1)}{2u}$ berechnet, einem Fehler von $\frac{3(uu+1)e\omega}{2u}$ bei der Form $\frac{1}{2}\lambda eu-\frac{\lambda e}{2u}$ oder einem Fehler von $3e\omega\operatorname{tang}F$ bei der Form $\lambda e\operatorname{tang}F$, sofern wir den Fehler in $\log\lambda$ oder $\log\frac{1}{2}\lambda$ vernachlässigen. Im ersten Fall lässt sich der Fehler auch durch
\SourceLeaf{050}$5e\omega\operatorname{tang}F$, im zweiten durch $\frac{3e\omega}{\cos F}$ ausdrücken. Daraus geht hervor, dass der Fehler im dritten Fall stets am kleinsten ist. Im ersten beziehungsweise zweiten Fall ist er größer, je nachdem, ob $u$ oder $\frac{1}{u}>2$ beziehungsweise $<2$ ist, oder je nachdem, ob $\pm F>36^0 52^{\prime}$ beziehungsweise $<36^0 52^{\prime}$ ist. — Der zweite Teil von $N$ ist dagegen stets einem Fehler von $\omega$ unterworfen.

VII. Umgekehrt ist klar, dass bei der Bestimmung von $u$ oder $F$ aus $N$ durch wiederholtes Probieren die Größe $u$ einem Fehler von $(1\pm5e\operatorname{tang}F)\omega\frac{\mathrm{d}u}{\mathrm{d}N}$ oder $\left(1+\frac{3e}{\cos F}\right)\omega\frac{\mathrm{d}u}{\mathrm{d}N}$ unterworfen ist, je nachdem, ob man den ersten Term im Ausdruck für $N$ in Faktoren oder in Teile zerlegt verwendet. Die Größe $F$ ist dagegen einem Fehler von $(1\pm3e\operatorname{tang}F)\omega\frac{\mathrm{d}F}{\mathrm{d}N}$ unterworfen. Die oberen Vorzeichen gelten nach dem Perihel, die unteren vor dem Perihel. Setzt man hier für $\frac{\mathrm{d}u}{\mathrm{d}N}$ oder für $\frac{\mathrm{d}F}{\mathrm{d}N}$ den Ausdruck $\frac{\mathrm{d}v}{\mathrm{d}N}$ ein, so erhält man die Wirkung dieses Fehlers auf die Bestimmung von $v$. Sie beträgt also $\frac{bb\operatorname{tang}\psi(1\pm5e\operatorname{tang}F)\omega}{\lambda rr}$ oder $\frac{bb\operatorname{tang}\psi(1+3e\operatorname{sec} F)\omega}{\lambda rr}$, wenn man die Hilfsgröße $u$ verwendet hat. Hat man dagegen $F$ verwendet, so ergibt sich für diese Wirkung
\[
\begin{aligned}
&=\frac{bb\operatorname{tang}\psi(1\pm3e\operatorname{tang}F)\omega}{\lambda rr}\\
&=\frac{\omega}{\lambda}\left\{\frac{(1+e\cos v)^2}{\operatorname{tang}\psi^3}\pm\frac{3e\sin v(1+e\cos v)}{\operatorname{tang}\psi^2}\right\}\raisebox{.45ex}{.}
\end{aligned}
\]
Soll der Fehler in Bogensekunden ausgedrückt werden, so ist der Faktor $206265^{\prime\prime}$ hinzuzufügen. Offenbar kann dieser Fehler nur dann beträchtlich werden, wenn $\psi$ ein kleiner Winkel ist, also wenn $e$ nur wenig größer als 1 ist. Für einige Werte von $e$ sind hier die größten Werte dieses dritten Ausdrucks bei Verwendung von sieben Dezimalstellen angegeben.

\begin{center}
\begin{tabular}{c|c}
$e$ & größter Fehler \\
1,3 & $0^{\prime\prime}{,}34$ \\
1,2 & $0{,}54$ \\
1,1 & $1{,}31$ \\
1,05 & $3{,}36$ \\
1,01 & $34{,}41$ \\
1,001 & $1064{,}65$ \\
\end{tabular}
\end{center}
\noindent Zu diesem aus einem fehlerhaften Wert von $F$ oder $u$ entstandenen Fehler muss man den in V bestimmten Fehler hinzurechnen, um die gesamte Unsicherheit von $v$ zu erhalten.

VIII. Löst man die Gleichung XI aus Artikel 22 mithilfe hyperbolischer Logarithmen und verwendet man dabei $F$ als Hilfsgröße, so erhält man durch ähnliche Überlegungen für die Wirkung des bei dieser Rechnung möglichen Fehlers auf die Bestimmung von $v$
\SourceLeaf{051}
\[
= \frac{(1+e\cos v)^2\,\omega'}{\mathrm{tang}\,\psi^3} \pm \frac{3e\sin v(1+e\cos v)\,\omega}{\lambda\,\mathrm{tang}\,\psi^2},
\]
wobei $\omega'$ die größte Unsicherheit in den Tafeln der hyperbolischen Logarithmen bezeichnet. Der zweite Teil dieses Ausdrucks stimmt mit dem zweiten Teil des in VII angegebenen Ausdrucks überein. Der erste Teil ist dagegen im Verhältnis $\lambda\omega' : \omega$ kleiner als der erste Teil des dortigen Ausdrucks, also im Verhältnis $1:23$, wenn man annehmen dürfte, dass die Tafel von \textsc{Ursinus} durchweg auf acht Stellen genau ist, also $\omega' = 0{,}000\,0000\,05$ gilt.
